\origpage{42}{ROBERTSON AND SEYMOUR}
\hspace*{12pt}(2) \textit{For each $(A, \Omega) \in \mathscr{S}$ either $A \subseteq H$ or $A$ has at most one circuit.}

\textit{Subproof.} Suppose first that $H \cap A$ has a circuit $C$. As in the proof of step (1) of (6.5), it follows that $|\bar{\Omega}| = 3$ and there are three disjoint paths of $H \cap A$ from $V(C)$ to $|\bar{\Omega}|$. But then by (1) (setting $A' = H \cap A$) we deduce that $A \subseteq H$ as required. We may assume then that $H \cap A$ is a forest. Let $B$ be a maximal forest of $A$ with $H \cap A \subseteq B$. If $E(B) = \emptyset$ then every edge of $A$ is a loop, and by (1), $|E(A)| \le 1$ as required. We may assume then that $E(B) \ne \emptyset$. By (1) (setting $A' = B$) we may assume that there exist $u, v \in \bar{\Omega}$ joined by a path of $A$ with no internal vertex in $\bar{\Omega}$ but not joined by any such path in $B$. Since $B$ is a maximal forest, there is a path $P$ of $B$ between $u$ and $v$, but it has an internal vertex $w \in \bar{\Omega}$. Then $\bar{\Omega} = \{u, v, w\}$. Let $B = B_1 \cup B_2$ where $V(B_1 \cap B_2) = \{w\}$ and $E(B_1 \cap B_2) = \emptyset$. Since there is a path of $A$ between $u$ and $v$ in $\delm{B}\fix{A} \setminus w$, there is an edge $e$ of $A$ with one end in $B_1 \setminus w$ and the other in $B_2 \setminus w$.\fixnote{fix:8.4-claim2}{First, ``in $B \setminus w$'' is a misprint for ``in $A \setminus w$''. Second, there is a gap: $B \setminus w$ may have more than two components, and then an edge $e$ between $B_1 \setminus w$ and $B_2 \setminus w$ need not join $u$ and $v$ in $(B + e) \setminus w$. For example, if $w$ has three neighbours in $B$, with components $T_u, T_v, T_3$ of $B \setminus w$ and $B_2 = T_v \cup T_3$, the edge $e$ may join $T_u$ and $T_3$. The proof can be completed by choosing $B$ with more care. Among the maximal forests of $A$ that contain $H \cap A$ and $P$, take $B$ such that the component $T_u$ of $B \setminus w$ that contains $u$ is as large as possible. Suppose that an edge $g$ of $A$ joins $T_u$ to a component $T$ of $B \setminus w$ that contains neither $u$ nor $v$. Then $T$ has no vertex in $\bar{\Omega}$. The forest $B$ has exactly one edge $f$ between $T$ and $w$: at least one, since otherwise $B + g$ would be a larger forest, and at most one, since $B$ has no circuit. If $f \in E(H)$, then the edges of $H \cap A$ inside $T$ or between $T$ and $w$ form a forest with at least one edge, and it has a leaf in $T$. This vertex is a leaf of $H \cap A$ not in $\bar{\Omega}$, which the rigidity of $H$ does not allow. So $f \notin E(H)$, and replacing $f$ by $g$ in $B$ gives a forest of the same kind with a larger component of $u$, which is impossible. Hence every edge of $A$ that leaves $T_u$, other than at $w$, ends in the component $T_v$ of $B \setminus w$ that contains $v$. A path of $A \setminus w$ from $u$ to $v$ leaves $T_u$, so some edge $e$ of $A$ joins $T_u$ to $T_v$, and $B + e$ joins $u$ and $v$ by a path that avoids $w$. So Claim (2) holds as printed.} By (1) (setting $A'$ to be the graph obtained by adding $e$ to $B$) we deduce that $A \setminus e = B$ and so $A$ has only one circuit. This proves (2).

Since every graph with at most one circuit can be drawn in the plane with all its vertices incident with the infinite region, we deduce from (2) and the rigidity of $H$ that for each $(A, \Omega) \in \mathscr{S}$, the drawing $H \cap A$ in the disc $\alpha(A, \Omega)$ can be extended to an embedding of $A$ in this disc, in such a way that each $v \in \bar{\Omega}$ is represented by the point $\alpha(v)$. It follows that we may assume, to simplify notation, that $G_1$ is a drawing in $\Sigma$.

Let $\Sigma'$ be the closure of the component of $\Sigma - U(C_1)$ which includes $U(C_2)$. Then $\Sigma'$ is a closed cylinder; let $A = G_1 \cap \Sigma'$, and let $B$ be the unique subgraph of $G_1$ such that $(A, B)$ is a separation and $V(A \cap B) = V(C_1)$. Then $(G, \Omega_1)$ is the composition of $(A \cup G_2, \Omega_0)$ with $(B, \Omega_1, \Omega_0)$, and the latter is ``rural'' in the sense of [3]. By [3, theorem (9.1)], we deduce that since $(G, \Omega_1)$ has a crooked transaction of order $\ge \pi$, so does $(A \cup G_2, \Omega_0)$. Now $(A \cup G_2, \Omega_0)$ is the composition of $(G_2, \Omega_2)$ and $(A, \Omega_0, \Omega_2)$, and the latter is a rural neighbourhood of breadth $\ge \pi + 5$, because $C_1, C_2$ have regional distance $\ge \pi + 4$. Because\fixnote{fix:8.4-proof}{A gap in the proof: (8.2) needs paths from $X$ to the inner boundary $\bar{\Omega}_2$ of the neighbourhood $(A, \Omega_0, \Omega_2)$, and here $\Omega_2 = \alpha^{-1}(F_2)$ comes from the end $F_2$ of the cylinder. The $2\pi + 8$ paths stop at $V(C_2)$, and in general they do not reach $\bar{\Omega}_2$, so (8.2) does not apply as written. The proof can be completed by cutting also at $C_2$, and applying (8.2) to the part of $A$ between $U(C_1)$ and $U(C_2)$, with a natural order of $V(C_2)$ in place of $\Omega_2$. This part is a rural neighbourhood of breadth at least $\pi + 5$, because $C_1, C_2$ have regional distance at least $\pi + 4$, and the $2\pi + 8$ paths run from $X$ to its inner boundary $V(C_2)$. The rest of the proof is unchanged, so (8.4) holds as printed.} there are $2\pi + 8$ disjoint paths of $A$ from $X$ to $V(C_2)$ with no vertex in $V(C_2)$ except their end vertices, it follows from (8.2) that $(A \cup G_2, \Omega_0)$ has a crooked transaction $\mathscr{P}$ of order $\ge \pi$ with $V(\mathscr{P}) \subseteq X$, such that no member of $\mathscr{P}$ has any internal vertex in $\bar{\Omega}_0$. Then $\mathscr{P}$ is a crooked transaction in $(G, \Omega_0)$; and since for each $P \in \mathscr{P}$,
\begin{equation*}
P \cap H \subseteq (A \cup G_2) \cap H = A \cap H = H \cap \Sigma'
\end{equation*}
and $U(P) \cap U(C_1) = \{u, v\}$ where $P$ has ends $uv$, it follows that $\mathscr{P}$ satisfies the theorem. $\square$
